% GENERATED FILE. Edit canonical lessons, then run npm run generate:books.
% canonical-source-hash: fnv1a64:a0e7d459e2e57fcd
% canonical-lessons: HI-C474-aql, HI-C474-bevkuf, HI-C474-nadan, HI-C474-pagal, HI-C474-bahakna

\chapter{Sense and Folly}
\label{ch:hi-sense-and-folly}
\hlchaptermodality{474}

\begin{chapteropening}
\textbf{By the end of this chapter:} \emph{I can say whether someone acted wisely or foolishly.}

Complete the last lesson of chapter 474: I can say whether someone acted wisely or foolishly.
\end{chapteropening}

\section[\texorpdfstring{\hi{अक़्ल} (aql)}{अक़्ल (aql)}]{\hi{अक़्ल} (aql) — sense; wits}

\label{lesson:HI-C474-aql}



\begin{quote}
Before the new one: say the Hindi for to grumble, then the Hindi for silence.
\end{quote}

\subsection*{You'll want to know: \hi{अक़्ल}}
\textbf{\hi{अक़्ल}} — \emph{aql} — \textquotedblleft{}sense; wits\textquotedblright{}.

\begin{quote}
\hi{अरुण} \hi{छोटा} \hi{है}, \hi{लेकिन} \hi{उसकी} \hi{अक़्ल} \hi{बहुत} \hi{तेज़} \hi{है।}
\end{quote}

\begin{quote}
\emph{Aruṇ choṭā hai, lekin uskī aql bahut tez hai.}
\end{quote}

\begin{quote}
Arun is young, but he's very sharp-witted.
\end{quote}

Feminine.

\begin{cousinweb}
\textbf{\hi{अक़्ल}} (\emph{aql}) is usually traced to Arabic \emph{ʿaql}, \textquotedblleft{}reason, intelligence\textquotedblright{}, which reached Hindi through Persian.
\end{cousinweb}

\begin{grammarlens}[title={What you've built}]
One more word for describing what you lived and felt.
\end{grammarlens}

\subsection*{Your turn}
\begin{itemize}
\raggedright
  \item \emph{Say it:} \emph{aql}
  \item \emph{Say it:} \emph{aql}, once more
  \item \emph{Say it:} \emph{Aruṇ choṭā hai, lekin uskī aql bahut tez hai.} — Arun is young, but he's very sharp-witted.
  \item \emph{Recall:} say the Hindi for to grumble, then the Hindi for silence, then say \emph{aql} again
\end{itemize}

\begin{tcolorbox}[breakable,colback=teal!4,colframe=teal!35!black,title={Before you move on}]
What does \hi{अक़्ल} mean? (\textquotedblleft{}Sense; wits\textquotedblright{}.) Say it once more.
\end{tcolorbox}
\section[\texorpdfstring{\hi{बेवकूफ़} (bevkūf)}{बेवकूफ़ (bevkūf)}]{\hi{बेवकूफ़} (bevkūf) — a fool; foolish}

\label{lesson:HI-C474-bevkuf}



\begin{quote}
Before the new one: say the Hindi for silence, then the Hindi for sense.
\end{quote}

\subsection*{You'll want to know: \hi{बेवकूफ़}}
\textbf{\hi{बेवकूफ़}} — \emph{bevkūf} — \textquotedblleft{}a fool; foolish\textquotedblright{}.

\begin{quote}
\hi{मैं} \hi{कितना} \hi{बेवकूफ़} \hi{हूँ}, \hi{छाता} \hi{घर} \hi{पर} \hi{भूल} \hi{गया}!
\end{quote}

\begin{quote}
\emph{maiṁ kitnā bevkūf hūṁ, chātā ghar par bhūl gayā!}
\end{quote}

\begin{quote}
What a fool I am — I left my umbrella at home!
\end{quote}

Masculine or feminine, by the person meant.

\begin{cousinweb}
\textbf{\hi{बेवकूफ़}} (\emph{bevkūf}) is usually traced to Persian \emph{bēvuqūf}, \textquotedblleft{}ignorant\textquotedblright{}, which joins Persian \emph{bē-}, \textquotedblleft{}without\textquotedblright{}, to Arabic \emph{vuqūf}, \textquotedblleft{}knowledge, understanding\textquotedblright{}.
\end{cousinweb}

\begin{grammarlens}[title={What you've built}]
One more word for describing what you lived and felt.
\end{grammarlens}

\subsection*{Your turn}
\begin{itemize}
\raggedright
  \item \emph{Say it:} \emph{bevkūf}
  \item \emph{Say it:} \emph{bevkūf}, once more
  \item \emph{Say it:} \emph{maiṁ kitnā bevkūf hūṁ, chātā ghar par bhūl gayā!} — What a fool I am — I left my umbrella at home!
  \item \emph{Recall:} say the Hindi for silence, then the Hindi for sense, then say \emph{bevkūf} again
\end{itemize}

\begin{tcolorbox}[breakable,colback=teal!4,colframe=teal!35!black,title={Before you move on}]
What does \hi{बेवकूफ़} mean? (\textquotedblleft{}A fool; foolish\textquotedblright{}.) Say it once more.
\end{tcolorbox}
\section[\texorpdfstring{\hi{नादान} (nādān)}{नादान (nādān)}]{\hi{नादान} (nādān) — naive; innocent}

\label{lesson:HI-C474-nadan}



\begin{quote}
Before the new one: say the Hindi for sense, then the Hindi for a fool.
\end{quote}

\subsection*{You'll want to know: \hi{नादान}}
\textbf{\hi{नादान}} — \emph{nādān} — \textquotedblleft{}naive; innocent\textquotedblright{}.

\begin{quote}
\hi{वह} \hi{बच्चा} \hi{अभी} \hi{नादान} \hi{है}, \hi{ख़तरा} \hi{नहीं} \hi{समझता।}
\end{quote}

\begin{quote}
\emph{vah baccā abhī nādān hai, khatrā nahīṃ samajhtā.}
\end{quote}

\begin{quote}
That child is still naive; he doesn't understand danger.
\end{quote}

\begin{cousinweb}
\textbf{\hi{नादान}} (\emph{nādān}) is usually traced to Persian \emph{nādān}, \textquotedblleft{}not knowing\textquotedblright{}, built on \emph{nā-}, \textquotedblleft{}not\textquotedblright{}, and \emph{dān}, \textquotedblleft{}knowing\textquotedblright{}. Hindi uses it kindly, for someone too young to know.
\end{cousinweb}

\begin{grammarlens}[title={What you've built}]
One more word for describing what you lived and felt.
\end{grammarlens}

\subsection*{Your turn}
\begin{itemize}
\raggedright
  \item \emph{Say it:} \emph{nādān}
  \item \emph{Say it:} \emph{nādān}, once more
  \item \emph{Say it:} \emph{vah baccā abhī nādān hai, khatrā nahīṃ samajhtā.} — That child is still naive; he doesn't understand danger.
  \item \emph{Recall:} say the Hindi for sense, then the Hindi for a fool, then say \emph{nādān} again
\end{itemize}

\begin{tcolorbox}[breakable,colback=teal!4,colframe=teal!35!black,title={Before you move on}]
What does \hi{नादान} mean? (\textquotedblleft{}Naive; innocent\textquotedblright{}.) Say it once more.
\end{tcolorbox}
\section[\texorpdfstring{\hi{पागल} (pāgal)}{पागल (pāgal)}]{\hi{पागल} (pāgal) — mad; crazy}

\label{lesson:HI-C474-pagal}



\begin{quote}
Before the new one: say the Hindi for a fool, then the Hindi for naive.
\end{quote}

\subsection*{You'll want to know: \hi{पागल}}
\textbf{\hi{पागल}} — \emph{pāgal} — \textquotedblleft{}mad; crazy\textquotedblright{}.

\begin{quote}
\hi{क्या} \hi{तुम} \hi{पागल} \hi{हो} \hi{गए} \hi{हो}?
\end{quote}

\begin{quote}
\emph{kyā tum pāgal ho gae ho?}
\end{quote}

\begin{quote}
Have you gone mad?
\end{quote}

\begin{cousinweb}
\textbf{\hi{पागल}} (\emph{pāgal}) is a tadbhava: CDIAL lists Hindi \emph{pāgal}, \textquotedblleft{}mad\textquotedblright{}, under \emph{*paggala}, a form it reconstructs (CDIAL 7643).
\end{cousinweb}

\begin{grammarlens}[title={What you've built}]
One more word for describing what you lived and felt.
\end{grammarlens}

\subsection*{Your turn}
\begin{itemize}
\raggedright
  \item \emph{Say it:} \emph{pāgal}
  \item \emph{Say it:} \emph{pāgal}, once more
  \item \emph{Say it:} \emph{kyā tum pāgal ho gae ho?} — Have you gone mad?
  \item \emph{Recall:} say the Hindi for a fool, then the Hindi for naive, then say \emph{pāgal} again
\end{itemize}

\begin{tcolorbox}[breakable,colback=teal!4,colframe=teal!35!black,title={Before you move on}]
What does \hi{पागल} mean? (\textquotedblleft{}Mad; crazy\textquotedblright{}.) Say it once more.
\end{tcolorbox}
\section[\texorpdfstring{\hi{बहकना} (bahaknā)}{बहकना (bahaknā)}]{\hi{बहकना} (bahaknā) — to go astray}

\label{lesson:HI-C474-bahakna}



\begin{quote}
Before the new one: say the Hindi for naive, then the Hindi for mad.
\end{quote}

\subsection*{You'll want to know: \hi{बहकना}}
\textbf{\hi{बहकना}} — \emph{bahaknā} — \textquotedblleft{}to go astray\textquotedblright{}.

\begin{quote}
\hi{बुरे} \hi{दोस्तों} \hi{के} \hi{साथ} \hi{रहकर} \hi{अरुण} \hi{का} \hi{भाई} \hi{बहक} \hi{गया।}
\end{quote}

\begin{quote}
\emph{bure dostoṁ ke sāth rahkar Aruṇ kā bhāī bahak gayā.}
\end{quote}

\begin{quote}
Mixing with bad friends, Arun's brother went astray.
\end{quote}

Past (he …): \hi{वह} \hi{बहक} \hi{गया}, \emph{vah bahak gayā}.

\begin{cousinweb}
\textbf{\hi{बहकना}} (\emph{bahaknā}) is inherited: CDIAL lists Hindi \emph{bahaknā}, \textquotedblleft{}to be misled\textquotedblright{}, under Sanskrit \emph{vyáthatē}, \textquotedblleft{}wavers, is disturbed\textquotedblright{} (CDIAL 12164).
\end{cousinweb}

\begin{grammarlens}[title={What you've built}]
One more word for describing what you lived and felt.
\end{grammarlens}

\subsection*{Your turn}
\begin{itemize}
\raggedright
  \item \emph{Say it:} \emph{bahaknā}
  \item \emph{Say it:} \emph{bahaknā}, once more
  \item \emph{Say it:} \emph{bure dostoṁ ke sāth rahkar Aruṇ kā bhāī bahak gayā.} — Mixing with bad friends, Arun's brother went astray.
  \item \emph{Recall:} say the Hindi for naive, then the Hindi for mad, then say \emph{bahaknā} again
\end{itemize}

\begin{tcolorbox}[breakable,colback=teal!4,colframe=teal!35!black,title={Before you move on}]
What does \hi{बहकना} mean? (\textquotedblleft{}To go astray\textquotedblright{}.) Say it once more.
\end{tcolorbox}
